\section{Fitting the polynomial}

The modelling stage takes the \texttt{Result} rows the calculator handed on and makes, for each
membrane separately, three fitted equations and one search. Nothing here looks at a reading or
a sample again. The whole stage is \texttt{src/limodel/analyzer/model/}, and
\texttt{Surrogate.find\_models} is the only entry point.

Two things are decided before any arithmetic. First, the membrane is not a variable: the runs
are grouped by \texttt{cem} and the equation is fitted once per membrane.
\texttt{docs/polynomial.md} allowed either that or a 0/1 indicator; the code took the
per-membrane branch, which means the fits never borrow strength from each other. Second, a
membrane needs at least \texttt{Surrogate.FEWEST\_RUNS}, which is 3, because an intercept and
one slope is two coefficients and leave-one-out needs a row to spare.

On the 16 calculated runs that leaves two membranes fitted and two refused:

\begin{table}[htbp]
  \centering
  \caption{The four membranes in the current data, and what the fitting did with each.}
  \label{tab:membranes-fitted}
  \begin{tabular}{@{}lrl@{}}
    \toprule
    Membrane & Runs & Outcome \\
    \midrule
    GO-10            & 9 & fitted in both settings \\
    Pristine         & 5 & fitted in current density; feed ratio dropped \\
    PDA/PEI/15CE/BT  & 1 & not fitted: 1 run against the 3 a fit needs \\
    PDA-GO10/PEI/BT  & 1 & not fitted: 1 run against the 3 a fit needs \\
    \bottomrule
  \end{tabular}
\end{table}

A membrane that cannot be fitted is not silently skipped. It goes into
\texttt{SurrogateModels\allowbreak{}.not\_fitted} with the reason attached, and the report prints
it, because
``this membrane has one run'' is the next experiment worth booking rather than a footnote.

\subsection{The two design columns, logged and coded}

\texttt{DesignPoints} owns both settings of every run and the two translations between real
units and the units the fitting works in. The columns are
\texttt{current\_density\_ma\_cm2} and \texttt{feed\_mg\_li\_ratio} --- current density $J$ in
$\unit{mA/cm^2}$, and the feed's Mg/Li in $\unit{mol/mol}$. Both come from the calculate stage:
$J$ is the median supply current over the membrane area, and the feed ratio is worked out from
the run's own feed concentrations.

Each setting is logged first, then centred and scaled:

\begin{align}
  \label{eq:coding}
  x_1 &= \frac{\ln J - c_J}{h_J}, &
  c_J &= \tfrac{1}{2}\bigl(\ln J_{\max} + \ln J_{\min}\bigr), &
  h_J &= \tfrac{1}{2}\bigl(\ln J_{\max} - \ln J_{\min}\bigr), \\
  \label{eq:coding-r}
  x_2 &= \frac{\ln R - c_R}{h_R}, &
  c_R &= \tfrac{1}{2}\bigl(\ln R_{\max} + \ln R_{\min}\bigr), &
  h_R &= \tfrac{1}{2}\bigl(\ln R_{\max} - \ln R_{\min}\bigr).
\end{align}

The extremes in \eqref{eq:coding} and \eqref{eq:coding-r} are taken over the runs of the
membrane being fitted, so $x_1$ and $x_2$ run from $-1$ to $+1$ by construction. That is the
whole point of the coding, and the search later leans on it.

Two reasons for the two transforms, and they are separate reasons.

The logarithm is because each setting spans several fold rather than a few units. Across the
GO-10 runs $J$ goes from $4.50$ to $25.67\unit{mA/cm^2}$, a factor of $5.70$, and the feed ratio
from $1.43$ to $11.42\unit{mol/mol}$, a factor of $8.00$. On a raw scale a quadratic over that
span is dominated by whatever happens at the top end; on a log scale the steps between the
levels run are roughly even, so every run gets comparable say in the curvature.

The centring and scaling is because the squared terms would otherwise swamp the linear ones.
Logged but uncoded, the GO-10 current density column runs over $1.504$ to $3.245$, so its
square would run over $2.26$ to $10.53$ --- a column an order of magnitude larger than its own
linear partner, strongly correlated with it, and pushing the least-squares solve towards a
badly conditioned matrix. In coded units $x$ and $x^2$
both sit inside $[-1,1]$ and the coefficients are directly comparable: $b_1$ and $b_3$ are both
``change in the outcome across half the measured span''.

Table~\ref{tab:coded} is the coding for the runs as they stand.

\begin{table}[htbp]
  \centering
  \caption{The GO-10 design, in real units and coded. Centres and half-widths are
           $c_J = 2.374635$, $h_J = 0.870558$, $c_R = 1.395778$, $h_R = 1.039721$ in logged
           units; $e^{c_J} = 10.75\unit{mA/cm^2}$ and $e^{c_R} = 4.04\unit{mol/mol}$ are the
           geometric centres.}
  \label{tab:coded}
  \begin{tabular}{@{}lrrrr@{}}
    \toprule
    Run & $J$ ($\unit{mA/cm^2}$) & $R$ ($\unit{mol/mol}$) & $x_1$ & $x_2$ \\
    \midrule
    AE & 13.67 &  1.43 & $+0.276058$ & $-1.000000$ \\
    AF & 13.67 &  2.86 & $+0.276058$ & $-0.333333$ \\
    AG & 13.67 &  5.71 & $+0.276058$ & $+0.333333$ \\
    AH &  7.50 &  1.43 & $-0.413220$ & $-1.000000$ \\
    AI &  7.50 &  1.43 & $-0.413220$ & $-1.000000$ \\
    AJ &  4.50 &  1.43 & $-1.000000$ & $-1.000000$ \\
    AN & 25.67 &  1.43 & $+1.000000$ & $-1.000000$ \\
    AP & 25.67 &  1.43 & $+1.000000$ & $-1.000000$ \\
    AT & 13.67 & 11.42 & $+0.276058$ & $+1.000000$ \\
    \bottomrule
  \end{tabular}
\end{table}

The five Pristine runs --- AK, AL, AM, AO, AQ --- sit at $J = 13.67$, $7.50$, $4.50$, $25.67$
and $25.67\unit{mA/cm^2}$, all at $R = 1.43\unit{mol/mol}$. Their $x_1$ values are the same four
numbers as GO-10's, because the extremes of $J$ happen to be the same, and their $x_2$ is $0.00$
throughout.

The shape of Table~\ref{tab:coded} is worth a moment, because it governs everything after it.
Current density was moved only at the lowest feed ratio, and the feed ratio was moved only at
$13.67\unit{mA/cm^2}$. The design is a cross, not a grid.

The code's columns differ in two places from the table in the agreed
document, and neither difference is a code fault:

\begin{itemize}[nosep]
  \item The doc lists $R$ levels of $5, 10, 20, 40, 60$. Those are mass ratios --- the feed
        holds $400\unit{mg/L}$ Li against $2000$, $4000$, $8000$ or $16000\unit{mg/L}$ Mg. The
        code's column is $\unit{mol/mol}$, so the same four levels read as $1.43$, $2.86$,
        $5.71$ and $11.42$, a fixed factor of $M_{\mathrm{Li}}/M_{\mathrm{Mg}} = 0.285538$ apart.
        Because \eqref{eq:coding-r} logs then centres, a constant factor moves $c_R$ and nothing
        else --- the coded values are identical either way. The $60$ level has not been run.
  \item The doc lists $J$ levels of $3, 6, 12, 24$. Those are nominal. The code uses the median
        supply current actually logged, which gives $4.50$, $7.50$, $13.67$ and
        $25.67\unit{mA/cm^2}$. The fit uses what the supply held, not what it was asked for.
\end{itemize}

\subsection{A setting held at one value is dropped, not a refusal}

The feed ratio never moved across the Pristine runs. A naive fit would then have a column of
identical numbers, which after centring is a column of zeroes --- the solve would be rank
deficient and the whole membrane would be thrown out.

\texttt{DesignPoints} handles it by name instead. \texttt{find\_scales} gives that setting a
half-width of exactly $0$; \texttt{find\_kept\_settings} returns only the settings whose
half-width is non-zero; \texttt{find\_usable\_coded\_names} turns those into the coded names the
polynomial knows, and \texttt{find\_dropped\_settings} names the rest so the report can say so.
A dropped setting is coded as $0.0$, the one place it was ever run, and its terms are left out by
name, so the zero column never reaches a least-squares call.

Today this applies to Pristine and only Pristine. What it changes:

\begin{itemize}[nosep]
  \item The candidate terms shrink from four to two. A term is kept only if the setting
        under it varies, which \texttt{find\_candidates\_on} decides, so Pristine's
        candidates are \texttt{x1} and \texttt{x1\_squared}.
  \item The report says it out loud. The ``fitted in'' line reads \emph{current density, with
        the feed ratio held at one value across these runs and never tested}, which stops a fit
        in one setting being read as a fit in two.
  \item The search still returns a feed ratio. \texttt{find\_real\_from\_coded} multiplies the
        coded value by a half-width of zero, so whatever $x_2$ the grid offers reads back as
        $1.43\unit{mol/mol}$ --- the one value it was held at.
\end{itemize}

Both settings held at one value is a different matter: \texttt{find\_scales} then returns
\texttt{None}, and \texttt{Surrogate.find\_refusal\_for} refuses the membrane with \emph{no
setting that varies across its runs, so there is nothing to fit}. Runs that moved nothing carry
no information about anything.

\subsection{The candidate terms, and why there is no interaction}

\texttt{SecondOrderPolynomial.TERMS\_IN\_ORDER} is five terms, not six:

\begin{equation}
  \label{eq:polynomial}
  \hat{y} \;=\; b_0 \;+\; b_1 x_1 \;+\; b_2 x_2 \;+\; b_3 x_1^2 \;+\; b_4 x_2^2
\end{equation}

\texttt{docs/polynomial.md} writes a sixth term, $b_5 \, J \cdot R$, and calls it the term that
says whether the best current depends on how much Mg is in the feed. It is not in the code, and
the reason is arithmetic rather than preference: over these runs that column does not carry any
information the other three do not already carry.

Here is the proof rather than the assertion. Look again at Table~\ref{tab:coded}. Every run with
$x_2 \neq -1$ sits at $x_1 = 0.276058$, and every run with $x_1 \neq 0.276058$ sits at
$x_2 = -1$. Write $u = 0.276058056371$, the coded value of $13.67\unit{mA/cm^2}$. Then for every
one of the nine GO-10 rows,

\begin{equation}
  \label{eq:identity}
  x_1 x_2 \;=\; u \;-\; x_1 \;+\; u\,x_2 .
\end{equation}

Check it both ways. On the arm where $x_2 = -1$, the right side is $u - x_1 - u = -x_1$, which is
$x_1 x_2$. On the arm where $x_1 = u$, it is $u - u + u x_2 = u x_2$, which is again $x_1 x_2$.
Every row lies on one arm or the other, so \eqref{eq:identity} holds everywhere.

Regressing the $x_1 x_2$ column on $[\mathbf{1}, x_1, x_2]$ with \texttt{numpy.linalg.lstsq}
returns the coefficients $(+0.276058, -1.000000, +0.276058)$ and a largest absolute residual of
$6.7 \times 10^{-16}$ --- floating-point dust, not a near miss. Checking \eqref{eq:identity}
directly gives a largest error of $1.1 \times 10^{-16}$.

So the interaction is an exact linear combination of the intercept and the two linear terms, and
the matrix ranks say the same thing:

\begin{table}[htbp]
  \centering
  \caption{Rank of the design matrix over the fitted runs, from
           \texttt{numpy.linalg.matrix\_rank}.}
  \label{tab:rank}
  \begin{tabular}{@{}lrrr@{}}
    \toprule
    Columns & Rows & Columns & Rank \\
    \midrule
    \multicolumn{4}{@{}l}{\textit{GO-10, 9 runs}} \\
    $[\mathbf{1},\,x_1,\,x_2,\,x_1 x_2]$                     & 9 & 4 & 3 \\
    $[\mathbf{1},\,x_1,\,x_2,\,x_1^2,\,x_2^2]$               & 9 & 5 & 5 \\
    $[\mathbf{1},\,x_1,\,x_2,\,x_1^2,\,x_2^2,\,x_1 x_2]$     & 9 & 6 & 5 \\
    \midrule
    \multicolumn{4}{@{}l}{\textit{Pristine, 5 runs}} \\
    $[\mathbf{1},\,x_1,\,x_2,\,x_1 x_2]$                     & 5 & 4 & 2 \\
    $[\mathbf{1},\,x_1,\,x_2,\,x_1^2,\,x_2^2,\,x_1 x_2]$     & 5 & 6 & 3 \\
    \bottomrule
  \end{tabular}
\end{table}

The singular values of the full six-column GO-10 matrix make the same point:

\begin{equation}
  \label{eq:singular}
  \sigma \;=\; \bigl(4.6488,\; 2.6124,\; 1.8651,\; 0.9843,\; 0.7760,\; 0.0000\bigr)
\end{equation}

The last entry of \eqref{eq:singular} is zero, not small. A zero singular value means the
interaction coefficient does not exist: infinitely many values of $b_5$, each with matching
changes in $b_0$, $b_1$ and $b_2$, fit the nine rows exactly as well. Least squares would hand
back one of them, and it would mean nothing.

The five-term matrix in \eqref{eq:polynomial} is full rank, 5 of 5, so the squared terms are not
affected by this. It is the product that has no support.

Pristine is simpler still: $x_2 = 0$ throughout, so $x_1 x_2$ is literally a column of zeroes.

The code does not just omit the term --- it also refuses any fit that turns out rank deficient.
\texttt{SecondOrderPolynomial.find\_fit} reads the rank back from \texttt{lstsq} and returns
\texttt{None} when it is short of the number of terms, so a term set these rows cannot settle
can never be reported as if it had been settled.

Getting the interaction back needs runs, not code: one run at a high feed ratio and a different
current density would break \eqref{eq:identity} on its own. AT moved the feed ratio to
$11.42\unit{mol/mol}$ but stayed at $13.67\unit{mA/cm^2}$, so it sits on the same arm.

\subsection{Forward stepwise selection}

The agreed document, \texttt{docs/polynomial.md}, writes the full equation and fits it. The code
does not. It grows the equation one term at a time, in
\texttt{FittedOutcome.find\_best\_polynomial}:

\begin{enumerate}[nosep]
  \item Start with the intercept alone and record its leave-one-out error.
  \item Of the terms it is allowed to add, try each one and keep the one that lowers the error
        most.
  \item Stop when no remaining term lowers it.
\end{enumerate}

Two rules shape the search. A squared term waits for its own linear term:
\texttt{find\_terms\_it\_may\_add} will not offer \texttt{x2\_squared} until \texttt{x2} is in,
because a curvature with no slope under it cannot be read as anything. And ties go to the
simpler term --- candidates are scanned in the order \texttt{x1}, \texttt{x2},
\texttt{x1\_squared}, \texttt{x2\_squared}, and an error that merely equals the incumbent does
not displace it, nor does it keep the search going.

This is a departure from the agreed document, and the reason is that the runs cannot settle that
many coefficients. Nine GO-10 rows for a five-term equation leaves four degrees of freedom, and
two of the nine rows repeat two others --- so the design holds seven distinct points. Fitting
five coefficients to seven points produces an equation that passes close to everything and
predicts nothing. Pristine is worse: five runs, four distinct currents. Growing
the equation only while each new term pays for itself out of sample is how the fit is kept to the
size the data can support.

The honest caveat is that the reported error is the same quantity the search minimised, so it is
mildly optimistic as an estimate of error on a fresh run. The comparison in the next subsection
is what guards against reading too much into it.

Once the term set is chosen, the coefficients that get printed come from a final fit on
\emph{all} the rows: \texttt{find\_fit} refits the chosen terms over every row after the
search finishes. The leave-one-out folds are only ever used to score term sets, never to
produce the published coefficients.

Table~\ref{tab:stepwise} is the search as it actually ran, for all six fits.

\begin{table}[htbp]
  \centering
  \caption{Every step of the forward search, with the leave-one-out error each candidate
           reached. The chosen term at each step is marked; a dash means the term was not
           offered, either because it was already in or because its linear partner was not.}
  \label{tab:stepwise}
  \begin{tabular}{@{}lrrrr@{}}
    \toprule
    Fit so far & \texttt{x1} & \texttt{x2} & \texttt{x1\_squared} & \texttt{x2\_squared} \\
    \midrule
    \multicolumn{5}{@{}l}{\textit{GO-10, permselectivity --- intercept alone: 0.7678}} \\
    intercept                     &  0.8683 & \textbf{0.5611} & ---    & ---    \\
    intercept, \texttt{x2}        &  0.7584 & ---             & ---    & 0.6051 \\
    \midrule
    \multicolumn{5}{@{}l}{\textit{GO-10, recovery \% --- intercept alone: 15.0521}} \\
    intercept                     & 16.2443 & \textbf{7.8979} & ---    & ---    \\
    intercept, \texttt{x2}        &  9.9497 & ---             & ---    & \textbf{7.1218} \\
    $\dots$, \texttt{x2\_squared} &  9.4710 & ---             & ---    & ---    \\
    \midrule
    \multicolumn{5}{@{}l}{\textit{GO-10, volts --- intercept alone: 4.0593}} \\
    intercept                     & \textbf{1.3385} & 4.3523  & ---    & ---    \\
    intercept, \texttt{x1}        & ---     & 1.5419          & 1.4018 & ---    \\
    \midrule
    \multicolumn{5}{@{}l}{\textit{Pristine, permselectivity --- intercept alone: 0.0614}} \\
    intercept                     &  0.0741 & ---             & ---    & ---    \\
    \midrule
    \multicolumn{5}{@{}l}{\textit{Pristine, recovery \% --- intercept alone: 2.9138}} \\
    intercept                     &  3.6629 & ---             & ---    & ---    \\
    \midrule
    \multicolumn{5}{@{}l}{\textit{Pristine, volts --- intercept alone: 5.8947}} \\
    intercept                     & \textbf{1.0682} & ---     & ---    & ---    \\
    intercept, \texttt{x1}        & ---     & ---             & 2.9568 & ---    \\
    \bottomrule
  \end{tabular}
\end{table}

Three things in that table are worth reading twice. No fit ever reached three terms except
GO-10's recovery, which took \texttt{x2} and then \texttt{x2\_squared}. Current density never
earned a place in a permselectivity or recovery fit --- for GO-10 it made both errors
\emph{worse} than the intercept alone. And two Pristine fits kept nothing at all: the best
prediction of its permselectivity or recovery at any setting is the average of the five runs.

\subsection{Leave-one-out, not a hold-out set}

\texttt{LeaveOneOutError.find\_rmse} is the only error measure in the stage. Each row takes its
turn being held back; the polynomial is fitted on the rest and asked to predict the one it never
saw; the misses are squared, averaged and rooted:

\begin{equation}
  \label{eq:loo}
  \mathrm{RMSE}_{\text{LOO}} \;=\; \sqrt{\frac{1}{n}\sum_{i=1}^{n}
    \bigl(\hat{y}^{(-i)}(x_i) - y_i\bigr)^2}
\end{equation}

Every row in \eqref{eq:loo} is predicted by a fit that did not include it, so the number is out
of sample $n$ times over. It returns \texttt{None} rather than a number whenever a fold cannot
be fitted, whenever an outcome is missing, and whenever there are fewer than two rows --- and the
stepwise search treats \texttt{None} as ``this term cannot be judged'' and moves on, which is
why a term set the data cannot settle never wins by default.

A train/test split was rejected for one reason: with nine rows a hold-out set measures nothing.
Hold two GO-10 runs back and the error is an average of two squared misses; swap which two and
the number moves more than the differences between the models it is meant to be choosing
between. Worse, the seven rows left would not span the design --- hold back AT and the feed
ratio no longer reaches $11.42\unit{mol/mol}$, so the model being tested is not the model being
proposed. Leave-one-out spends no rows permanently and produces one honest out-of-sample miss
per run instead of two in total.

The real test set is arriving on its own. AV, AW, AX and AY have no lab results yet, so they are
dropped today; when their results land they are unseen runs that no fit has touched.

\subsection{Replicate scatter: the yardstick}

An error of $0.561$ on permselectivity is meaningless on its own. It needs something to be
compared against, and the only fair comparison is how far apart the bench's own repeats came
out. \texttt{ReplicateScatter} works that out.

A repeat is two runs at the same membrane and the same two settings, matched to two decimal
places --- \texttt{ReplicateScatter.PLACES}. The rounding is there because the settings are
computed, not typed in: $25.666666666666668$ and $25.666666666666664$ are the same experiment,
and matching the raw floats would miss it. Runs missing the outcome or either setting are passed
over.

Within each group the spread is pooled across groups rather than averaged, which is the standard
way to combine two small samples of the same underlying noise:

\begin{equation}
  \label{eq:pooled}
  s \;=\; \sqrt{\frac{\sum_{g}\sum_{i \in g}\bigl(y_i - \bar{y}_g\bigr)^2}
                     {\sum_{g}\bigl(n_g - 1\bigr)}}
\end{equation}

The denominator in \eqref{eq:pooled} is the degrees of freedom, and it is reported beside the
spread so a reader can see how thin the yardstick is. No group with a repeat in it means no
yardstick at all, and \texttt{find\_spread} returns \texttt{None}.

Three design points repeat in the current data:

\begin{table}[htbp]
  \centering
  \caption{The replicated design points and what each pair gave, with the pooled spread from
           \eqref{eq:pooled}. $J$ is in $\unit{mA/cm^2}$ and $R$ in $\unit{mol/mol}$. GO-10
           pools two pairs; Pristine has one.}
  \label{tab:replicates}
  \begin{tabular}{@{}lrrlll@{}}
    \toprule
    Runs & $J$ & $R$ & Permselectivity & Recovery (\%) & Volts \\
    \midrule
    AH, AI &  7.50 & 1.43 & 2.3881, 2.0322 & 36.59, 37.66 & 10.55, 10.22 \\
    AN, AP & 25.67 & 1.43 & 1.9177, 1.3217 & 39.50, 26.32 & 18.54, 16.91 \\
    \multicolumn{3}{@{}l}{GO-10 pooled, 2 d.o.f.}
           & 0.3471 & 6.6124 & 0.8315 \\
    \midrule
    AO, AQ & 25.67 & 1.43 & 0.2115, 0.2986 &  7.83, 11.38 & 16.84, 17.04 \\
    \multicolumn{3}{@{}l}{Pristine pooled, 1 d.o.f.}
           & 0.0615 & 2.5075 & 0.1414 \\
    \bottomrule
  \end{tabular}
\end{table}

Look at the GO-10 recovery pair AN and AP: $39.50\%$ and $26.32\%$ at the same setting on the
same membrane. That single pair is most of the $6.6124$ pooled spread, and it is the honest
statement of how reproducible recovery is at $25.67\unit{mA/cm^2}$ today.

This is why \texttt{FittedOutcome} carries the spread alongside the coefficients and the error:
the three travel together because an equation on its own cannot be judged. A fit that misses by
more than the bench's own repeats has learnt nothing the repeats did not already say --- it has
fitted the scatter. The report makes that call outright rather than printing two numbers and
leaving it to the reader, and it makes it at the three decimals it prints, so two numbers that
display the same are not treated as separated.

\subsection{The gated search for the best setup}

The question is not ``what is the most selective setting'' but ``what is the most selective
setting that is still worth running''. So all three fitted outcomes are needed:
permselectivity to maximise, and recovery and voltage to gate.
\texttt{BestSetupUnderTheGates} does this per membrane.

\begin{itemize}[nosep]
  \item \textbf{The grid.} \texttt{STEPS} is 21 levels per setting, evenly spaced in coded units
        from $-1$ to $+1$, which is $0.1$ apart --- $441$ combinations in all. Because the coding
        is on logs, those levels are geometric in real units: each step multiplies $J$ by
        $1.0910$ and the feed ratio by $1.1096$.
  \item \textbf{The sweep stays inside the measured range.} Coded $-1$ and $+1$ are exactly the
        lowest and highest run of the membrane being searched, so sweeping that span and no
        further is what keeps the answer to a setting the bench has evidence for. For GO-10 that
        is $4.50$ to $25.67\unit{mA/cm^2}$ and $1.43$ to $11.42\unit{mol/mol}$. A quadratic
        extrapolated past its last data point curves away fast and in whichever direction the
        algebra happens to point; the recommendation would be a reading taken past the last run.
  \item \textbf{The recovery floor.} \texttt{RECOVERY\_FLOOR\_PCT} is $35.0$. Less lithium than
        that recovered is not worth a run.
  \item \textbf{The voltage ceiling.} \texttt{VOLTAGE\_CEILING\_V} is $20.0$. The supply and the
        stack are not asked to go above it.
  \item \textbf{The winner.} Of the points that clear both gates, the one with the highest
        predicted permselectivity, converted back to $\unit{mA/cm^2}$ and $\unit{mol/mol}$ by
        \texttt{find\_real\_from\_coded}.
\end{itemize}

Both gates must be clear; \texttt{find\_gates\_failed\_at} returns the list of gates a point
falls foul of, and only an empty list survives. A point where any one of the three outcomes
cannot be predicted returns \texttt{None} and is counted as neither passed nor failed.

Refusing every point is a finding and not an error, and the code is built around that. It keeps a
count of how many of the $441$ points each gate turned away, and \texttt{find\_refusal} names the
gate responsible in a sentence the report prints. ``Nothing in the range you have tested is worth
running'' is a true and useful answer to the question that was asked.

That is exactly what happens today. For both fitted membranes, all $441$ points fail the recovery
floor and $0$ fail the voltage ceiling:

\begin{table}[htbp]
  \centering
  \caption{The gated search on the current fits. The highest predicted values anywhere on the
           grid, against the two limits.}
  \label{tab:gates}
  \begin{tabular}{@{}lrrrrr@{}}
    \toprule
    Membrane & Points & Best recovery (\%) & Floor (\%) & Worst volts & Ceiling \\
    \midrule
    GO-10    & 441 & 33.285 & 35.0 & 18.246 & 20.0 \\
    Pristine & 441 &  9.057 & 35.0 & 17.320 & 20.0 \\
    \bottomrule
  \end{tabular}
\end{table}

GO-10 comes close --- $33.285\%$ against a floor of $35\%$ --- and the voltage was never the
binding constraint for either membrane. Pristine fails by a mile for a simple reason: its
recovery fit kept no terms, so it predicts $9.057\%$ everywhere, which is the mean of its five
runs.

\subsection{What the fits are today}

Table~\ref{tab:fits} is what \texttt{uv run limodel analyze} prints. The terms column is the
equation exactly as the report writes it; the error is the leave-one-out RMSE from
\eqref{eq:loo}; the scatter is the pooled replicate spread from \eqref{eq:pooled} with its
degrees of freedom.

\begin{table}[htbp]
  \centering
  \caption{Every fitted equation in the current data, with its leave-one-out error and the
           pooled bench scatter beside it; the scatter's degrees of freedom are in brackets.
           Only Pristine's permselectivity is inside the scatter.}
  \label{tab:fits}
  \begin{tabular}{@{}llrr@{}}
    \toprule
    Outcome & Terms as printed & Error & Scatter \\
    \midrule
    \multicolumn{4}{@{}l}{\textit{GO-10 --- fitted in current density and the feed ratio}} \\
    permselectivity & \texttt{intercept +1.126, x2 -0.712}      & 0.561 & 0.347 (2) \\
    recovery \%     & \texttt{intercept +11.58, x2 -15.63,}
                                                                 & 7.122 & 6.612 (2) \\
                    & \texttt{x2\_squared +6.08}                &       &           \\
    volts           & \texttt{intercept +12.7, x1 +5.545}       & 1.339 & 0.832 (2) \\
    \midrule
    \multicolumn{4}{@{}l}{\textit{Pristine --- fitted in current density, with the feed ratio
                                   held at one value}} \\
    permselectivity & \texttt{intercept +0.2326}                & 0.061 & 0.062 (1) \\
    recovery \%     & \texttt{intercept +9.057}                 & 2.914 & 2.508 (1) \\
    volts           & \texttt{intercept +11.41, x1 +5.905}      & 1.068 & 0.141 (1) \\
    \bottomrule
  \end{tabular}
\end{table}

Five of the six fits print the sentence \emph{which is worse than the bench repeating itself, so
this fit has earned no conclusion}. Neither membrane gets a best setup; both print \emph{nothing
inside the measured range cleared the lithium recovery floor}.

Read plainly, that is the state of the modelling. One fit --- Pristine permselectivity, at
$0.061$ against a scatter of $0.062$ --- is inside the bench's own noise, and it is the fit that
kept no terms at all, so what it says is that Pristine's permselectivity is about $0.23$
regardless of current density. Everything else misses the held-out run by more than two repeats
of the same run differ from each other.

The two coefficients that do point somewhere are worth naming even so, because they are the only
two the search kept more than once. GO-10's permselectivity falls with the feed ratio
($-0.712$ across half the coded span, roughly a factor of $2.83$ in Mg/Li), and voltage rises
with current density on both membranes by about the same amount ($+5.545$ and $+5.905$ volts
across half the span). Neither is a surprise, and neither has cleared the yardstick.

The useful conclusion is about the design rather than the chemistry. The cross in
Table~\ref{tab:coded} is why there is no interaction term, five of six errors are above the
scatter, and the two thinnest membranes have one run each. Runs at a high feed ratio and a
current density other than $13.67\unit{mA/cm^2}$ would break \eqref{eq:identity} and let the
interaction be fitted; a third run at $25.67\unit{mA/cm^2}$ on GO-10 would tighten the recovery
yardstick that AN and AP currently set at $6.6$ percentage points.
